% !PRINTED 280
\weberpage{280}
Wir können daher dasselbe Verfahren anwenden wie oben, wenn wir noch
die Beschränkung hinzufügen, daß nur \gesperrt{symmetrische} Funktionen
von $A$ und $B$, d. h. rationale Funktionen von $AB$, $A+B$ benutzt
werden, weil \gesperrt{nur} unter dieser Voraussetzung aus einem der Werte einer
% [sic] Die Quelle verweist auf (17); die einschlägige Vertauschungstafel ist (21), während (17) die vorhergehende Argumentliste ist.
solchen Funktion durch die Vertauschungen (17) alle übrigen folgen. Man
berechnet leicht die folgenden Beispiele:
\weq{22}{
\begin{array}{l}
n = 15, \quad AB+1 = 0, \\
n = 35, \quad 2(A+B)-AB = 5 \\
n = 39, \quad 2(A+B)-AB = 3.
\end{array}
}

Die Rechnung bietet auch in noch komplizierteren Fällen keine
unüberwindlichen Schwierigkeiten. So habe ich in der Abhandlung Acta
mathematica Bd. II, S. 359 für den Fall $n = 105$ folgende Formel
mitgeteilt:
\[
\begin{array}{l}
A = \sum \dfrac{f(3\omega)f(5\omega)f(7\omega)f(105\omega)}{f(\omega)f(15\omega)f(21\omega)f(35\omega)}, \\[2ex]
B = \sum \dfrac{f(\omega)f(15\omega)f(21\omega)f(35\omega)}{f(3\omega)f(5\omega)f(7\omega)f(105\omega)}, \\[2ex]
A^2+B^2-4(A+B)^3+10AB(A+B) \\
\quad{}+4(A+B)^2+10AB+14(A+B)+5 = 0,
\end{array}
\]
wenn sich die Summen $\Sigma$ in $A$ und $B$ auf die drei Funktionen
$f, f_1, f_2$ beziehen.

\begin{center}
§ 77.\quad \gesperrt{Geometrische Deutung der irrationalen Modulargleichungen als}\\
\gesperrt{Modularkorrespondenzen.}
\end{center}

Den Inhalt der irrationalen Formen der Modulargleichungen machen wir
durch eine geometrische Deutung anschaulicher. Betrachten wir zunächst
den Fall $n \equiv -1 \ (\mathrm{mod}\ 8)$ und setzen
\weq{1}{x = f(\omega), \quad y = f_1(\omega), \quad z = f_2(\omega),}
so wird hierdurch, wenn wir $x, y, z$ als Cartesische Koordinaten eines
Punktes $P$ ansehen, eine Raumkurve dargestellt, die wir auch durch die
beiden Gleichungen:
\weq{2}{
\begin{array}{c}
x^8 - y^8 - z^8 = 0 \\
xyz = \sqrt{2}
\end{array}
}
ausdrücken können, und die wir die \gesperrt{Grundkurve} nennen wollen.
Diese Kurve ist also von der 24. Ordnung. Man kann sie auf
%CONT>>
